\documentclass[11pt]{article}

\usepackage[margin=1in]{geometry}
\usepackage{booktabs}
\usepackage{graphicx}
\usepackage{amsmath}

\setlength{\emergencystretch}{2em}

\title{Weekly Report: Week 6}
\author{}
\date{2026-09-02}

\begin{document}
\maketitle

All results use the matcher
\texttt{deepseek-\allowbreak v4-\allowbreak flash-\allowbreak 0731} and the embedding
model \texttt{bge-\allowbreak large-\allowbreak en-\allowbreak v1.5}. Every comparison within a dataset is scored over an identical
candidate set: 2{,}665 pairs on Fodors--Zagat, 13{,}080 on DBLP--ACM, 88{,}296 on
Amazon--Walmart.

\section{Pipeline Overview}

The pipeline runs in five stages, each timed and recorded in a per-run JSON summary under
\texttt{logs/<dataset>/runs/}.

\textbf{Load and normalize.} Both tables and the ground truth are read, ID columns are
dropped, and every remaining field is normalized: HTML entity references are decoded, text
is Unicode-decomposed with combining marks stripped, and address and directional
abbreviations are standardized. This matters more than it appears: \texttt{ACM.csv} stores
461 HTML entity references (\texttt{\&\#246;}, \texttt{\&mdash;}) where \texttt{DBLP.csv}
stores none, so without decoding, the same author name is two unrelated strings across the
two sources.

\textbf{Embed and block.} Each attribute column is encoded independently with SentenceBERT,
the column vectors are concatenated and L2-normalized into one embedding per record, and
multi-probe locality-sensitive hashing retrieves the top-$k$ most similar Table-B candidates
for every Table-A record. On DBLP--ACM this reduces 6{,}001{,}104 possible pairs to 13{,}080
candidates while retaining all 2{,}224 ground-truth matches.

\textbf{Step 1 --- global tuple condensation (optional).} One attribute subset is chosen
once for the whole dataset and applied to every record. The tables may optionally be
re-blocked on the condensed attributes; see Section~\ref{sec:reblock}.

\textbf{Step 1.5 --- pair-adaptive attribute selection (optional).} A small sample of
candidate pairs is profiled by the LLM for per-attribute importance, and a hybrid
predictor/retriever model transfers a \emph{different} attribute subset to every candidate
pair. Blocking is unchanged, so the candidate set is identical to the no-selection baseline.

\textbf{Match.} One LLM call decides Yes/No per candidate pair, using only the attributes
selected for that pair.

\textbf{Metrics and logging.} Precision, recall, and F1 are computed end-to-end, against all
ground-truth matches, so pairs lost at blocking count against recall.

\section{Methodology}
\label{sec:methodology}

\subsection{Step 1.5: Pair-Adaptive Attribute Selection}

Step~1.5 has a profiling stage and a transfer stage.

\textbf{Profiling.} A label-free sampler selects a small set of candidate pairs (20 on
Fodors--Zagat, 50 on DBLP--ACM) and the LLM is asked, for each, which attributes matter for
deciding whether the two records describe the same entity. Ground truth is not consulted at
any point. The \texttt{diversity} sampler is used throughout this report, holding the
sampler fixed so that other variables can be read.

\textbf{Transfer.} For every candidate pair, six cheap similarity features are computed per
attribute --- token overlap, edit similarity, exact match, semantic similarity, numeric
similarity, and a numeric-type indicator --- without calling the LLM. Two models then
predict per-attribute importance from those features: a Ridge regressor fit on the profiled
pairs, and a nearest-neighbour retriever that averages the importance vectors of the five
most feature-similar profiled pairs. Their outputs are blended equally and normalized to sum
to one, and a selection rule decides how many of the ranked attributes to keep.

The profiling stage costs about 1\% of a run's tokens. Its purpose is to buy a rule that can
be applied to the remaining 99\% without further LLM calls.

\subsection{Two Ways to Elicit Importance}
\label{sec:scoring}

The profiling prompt is one of the two variables this report isolates. Two formulations were
compared, identical in framing and guardrails and differing only in the answer requested.

\textbf{Ordinal.} The LLM assigns each attribute an integer score from 0 (not useful) to 3
(decisive), scoring each attribute independently.

\textbf{Decisive.} The LLM returns the smallest subset of attributes sufficient to make the
decision, weighing attributes against one another in a single judgement.

The motivation for the second formulation was that the recorded 0--3 scores were not being
used discriminatively. Aggregated over all runs, DBLP--ACM used score 0 on 1\% of
judgements, collapsing a four-point scale to three; Fodors--Zagat used score 3 on 55\%. An
importance vector that is nearly flat cannot express that one attribute matters more than
another, whatever selection rule consumes it.

\subsection{How Many Attributes to Keep}
\label{sec:policy}

Profiling and transfer produce a ranked importance vector per pair. A separate rule turns
that ranking into a subset, and it is the second variable isolated here.

\textbf{Cumulative} keeps attributes in descending importance until their shares sum past a
fixed target, by default 0.8. It is the pipeline's historical default and is not scale-free:
importance is normalized to sum to one, so on a wide schema each attribute's share shrinks
and a fixed target becomes unreachable. On Amazon--Walmart's 13 attributes the 0.8 target
needs 7.7 attributes on average against a cap of 5, so the cap binds for every pair and the
rule cannot express ``this pair needs fewer fields'' at all.

\textbf{Gap} cuts at the largest drop in the ranked importance profile. It has no threshold
to tune and does not saturate on any of the three schemas.

Both rules are subject to a floor and a cap on the number of attributes kept. The floor was
2 by default and is set to 1 wherever a run is allowed to reduce a pair to a single
attribute.

\subsection{Restricting Selection to Attributes Carrying Evidence}
\label{sec:mask}

An attribute both records leave blank supports no decision either way, but the transfer
model can still rank it first: its features are computed from two empty strings and the
profiling signal that trained the ranking came from pairs where the field was populated.
Selecting it hands the matcher a prompt whose only content is two blank fields.

Before any rule is applied, each pair is therefore given a mask marking the attributes
present on both sides. Masked attributes have their predicted importance set to zero, the
vector is renormalised over the remainder, and the ranking the rule reads contains only
usable attributes; if a pair has none, the unmasked ranking is used unchanged rather than
selecting nothing.

The correction is confined to wide, sparsely populated schemas. On Amazon--Walmart it binds
for 23{,}439 of 88{,}296 pairs (26.5\%) under decisive elicitation, and every affected pair
had been selecting a field blank on at least one side. On DBLP--ACM it binds for 5 pairs of
13{,}080 and leaves F1 unchanged at 0.9397, so the results of Section~\ref{sec:dblp} stand
as reported. Ordinal elicitation is unaffected on both datasets: it never ranked a blank
attribute first to begin with, and its selections are byte-identical before and after.

\subsection{Hyperparameters}
\label{sec:hyper}

Everything below is held fixed across every run reported here unless a result explicitly
varies it. The two quantities that were varied deliberately --- the elicitation format and
the selection rule --- are the subject of Sections~\ref{sec:dblp} and~\ref{sec:rule}.

\begin{table}[h]
\centering
\begin{tabular}{lll}
\toprule
Stage & Parameter & Value \\
\midrule
Embedding   & model                       & \texttt{BAAI/bge-large-en-v1.5} \\
            & encoding                    & per column, concatenated, L2-normalised \\
\midrule
Blocking    & scheme                      & multi-probe LSH, random hyperplanes \\
            & tables / planes / flips     & 15 / 8 / 1 (DBLP--ACM), 10 / 4 / 1 (Amazon--Walmart) \\
            & top-$k$ per Table-A record  & 5 (DBLP--ACM), 4 (Amazon--Walmart) \\
\midrule
Profiling   & sampler                     & \texttt{diversity} \\
            & sample size                 & 50 candidate pairs \\
            & elicitation                 & \texttt{decisive} (\texttt{ordinal} where stated) \\
            & concurrency                 & 4 workers \\
\midrule
Transfer    & predictor                   & Ridge regression over 6 features per attribute \\
            & retriever                   & cosine $k$-NN, $k = 5$ \\
            & fusion                      & \texttt{weighted\_average}, 0.5 / 0.5 \\
\midrule
Selection   & rule                        & \texttt{gap} (no threshold) \\
            & $k$ bounds                  & $\min = 1$, $\max = 5$ \\
            & cumulative target           & 0.8 (only for the \texttt{cumulative} rule) \\
            & ratio $\lambda$             & 1.5 (only for the \texttt{ratio} rule) \\
            & required roles              & \texttt{identity} \\
            & attribute order in prompt   & canonical (table-column order) \\
\midrule
Matcher     & model                       & \texttt{deepseek/deepseek-v4-flash-0731} \\
            & temperature                 & 0 \\
            & max tokens                  & 256 \\
            & reasoning                   & disabled \\
            & concurrency                 & 8 workers \\
\midrule
Seeds       & pipeline / blocking         & 42 / 42 \\
\bottomrule
\end{tabular}
\caption{Hyperparameters. Attribute \emph{roles} are assigned per dataset from the column
name: \texttt{title} and \texttt{modelno} are \texttt{identity}, \texttt{authors} and
\texttt{brand} \texttt{strong\_support}, \texttt{venue}, \texttt{year} and
\texttt{category} \texttt{context}, and the remaining Amazon--Walmart columns
\texttt{other}. Only \texttt{identity} is required, which is what forces at least one
identifying field into every selection.}
\end{table}

\section{Experimental Design}
\label{sec:design}

\subsection{Why Comparisons Are Paired}
\label{sec:paired}

The matcher LLM is not deterministic, and on these datasets its variation is large enough to
swamp the effects being measured. Running the no-selection baseline twice on Fodors--Zagat,
with identical prompts and identical attributes, produced F1 of 0.8917 and 0.9596 --- a gap
of 0.068, from 19 of 2{,}665 answers changing, 18 of them from Yes to No. Precision moved
from 0.823 to 0.947 while recall was unchanged.

The same measurement on DBLP--ACM gives a much smaller floor: two independent draws of the
no-selection baseline produced F1 of 0.8706 and 0.8613, a gap of 0.0093. This is expected ---
DBLP--ACM has 2{,}224 ground-truth matches against Fodors--Zagat's 110, so the same absolute
number of flipped answers perturbs the metrics far less.

Comparing two methods by their F1 values from separate runs therefore measures the matcher's
variance more than it measures the methods. Two consequences follow.

First, the matcher response cache is enabled for all comparisons. Where two methods select
the same attributes for a pair, the pair yields one shared answer and cannot contribute a
spurious difference. Section~\ref{sec:mechanism} shows this working: across the 9{,}803 pairs
where two very different methods happened to select the same single attribute, they recorded
identical false-positive counts.

Second, methods are compared by McNemar's test on the pairs where they actually differ,
rather than by comparing aggregate F1. This is a paired test on the same candidate set and
the same matcher draw, and it is the only comparison in this report that carries a
significance claim.

A caveat this design imposes: because repeated runs of one configuration replay the cache,
they are not independent measurements and their standard deviation is meaningless.
Replication across seeds was dropped in favour of the paired design.

\subsection{Holding the Candidate Set Fixed Across Methods}
\label{sec:reblock}

Step~1.5 leaves blocking untouched, so its candidate set is identical to the no-selection
baseline. Step~1 condenses the record and can optionally re-block on the condensed
attributes, which would change which pairs are compared. It is therefore run with
re-blocking disabled (\texttt{--step1-no-reblock}) throughout, so that every configuration is
scored over the same 13{,}080 candidate pairs containing the same 2{,}224 ground-truth
matches, and attribute choice is the only difference between them.

\section{Results}

\subsection{Fodors--Zagat is saturated}

\begin{table}[h]
\centering
\begin{tabular}{lrrrrrrr}
\toprule
Method & Attrs/pair & TP & FP & FN & F1 & Precision & Recall \\
\midrule
Step 1.5, ordinal      & 4.00 & 108 & 6 & 2 & \textbf{0.9643} & 0.9474 & \textbf{0.9818} \\
Step 1.5, decisive     & \textbf{2.68} & 106 & \textbf{4} & 4 & 0.9636 & \textbf{0.9636} & 0.9636 \\
No selection           & 5.00 & 107 & 6 & 3 & 0.9596 & 0.9469 & 0.9727 \\
Step 1, supervised     & 2.00 & 106 & 5 & 4 & 0.9593 & 0.9550 & 0.9636 \\
\bottomrule
\end{tabular}
\caption{Fodors--Zagat, 110 ground-truth matches, \texttt{class} excluded.}
\end{table}

Every pairwise McNemar test returns $p = 1.00$; the largest disagreement between any two
methods is 9 pairs out of 2{,}665. Doing nothing already recovers 107 of 110 matches, so
there is no headroom for a method to distinguish itself, and with 110 matches the matcher's
own instability moves F1 further than any method does. The defensible claim here is a 46\%
reduction in attributes at unchanged accuracy.

\subsection{DBLP--ACM separates the methods}
\label{sec:dblp}

\begin{table}[h]
\centering
\resizebox{\textwidth}{!}{%
\begin{tabular}{llrrrrrrr}
\toprule
Method & Rule & Attrs/pair & TP & FP & FN & F1 & Precision & Recall \\
\midrule
Step 1.5, decisive & gap            & 1.42 & 2197 & \textbf{255} & 27 & \textbf{0.9397} & \textbf{0.8960} & 0.9879 \\
Step 1, supervised (200 labels) & --- & 1.00 & 2198 & 339 & 26 & 0.9233 & 0.8664 & 0.9883 \\
Step 1.5, decisive & cumulative 0.8 & 2.24 & 2214 & 535 & 10 & 0.8904 & 0.8054 & 0.9955 \\
No selection       & ---            & 4.00 & 2216 & 651 &  8 & 0.8706 & 0.7729 & \textbf{0.9964} \\
Step 1.5, ordinal  & cumulative 0.8 & 3.77 & 2215 & 699 &  9 & 0.8622 & 0.7601 & 0.9960 \\
Step 1, supervised (10 labels) & --- & 2.00 & 2211 & 739 & 13 & 0.8547 & 0.7495 & 0.9942 \\
\bottomrule
\end{tabular}}
\caption{DBLP--ACM, 2{,}224 ground-truth matches.}
\end{table}

\begin{table}[h]
\centering
\begin{tabular}{lrrl}
\toprule
Comparison (A vs.\ B) & A right & B right & $p$ \\
\midrule
Decisive/gap vs.\ no selection            & 527 & 150 & $5.2\times10^{-50}$ \\
Decisive/gap vs.\ supervised (200)        & 112 &  29 & $1.0\times10^{-12}$ \\
Decisive/gap vs.\ decisive/cum.\ 0.8      & 375 & 112 & $3.5\times10^{-34}$ \\
Supervised (200) vs.\ decisive/cum.\ 0.8  & 389 & 209 & $1.6\times10^{-13}$ \\
Decisive/cum.\ 0.8 vs.\ no selection      & 468 & 350 & $4.2\times10^{-5}$ \\
Decisive/cum.\ 0.8 vs.\ ordinal           & 425 & 258 & $1.7\times10^{-10}$ \\
No selection vs.\ ordinal                 & 135 &  86 & $1.2\times10^{-3}$ \\
Ordinal vs.\ supervised (10)              & 497 & 453 & $0.16$ \\
\bottomrule
\end{tabular}
\caption{McNemar tests over the pairs on which two methods disagree.}
\end{table}

The best configuration reaches F1 0.9397 on 1.42 of 4 attributes per pair, beating both no
selection and a global baseline trained on 400 labelled pairs, while using no labels
itself. Every figure in this table is one draw of the matcher, recorded between 12 and 19
August; Section~\ref{sec:draws} repeats the two headline rows and shows that the absolute
values do not reproduce two weeks later while the ordering does.

Recall is close to constant across the table (0.988--0.996) while precision spans
0.750--0.896. Attribute selection does not help the matcher find matches it was missing; it
stops the matcher asserting matches that are not there. The one cost appears at the
aggressive end, where the configurations showing barely more than one attribute give up
18--19 of 2{,}224 matches.

Holding the rule fixed, the decisive formulation beats no selection
($p = 4.2\times10^{-5}$) while the ordinal formulation is significantly \emph{worse} than
no selection ($p = 1.2\times10^{-3}$). Only the profiling question differs between them.
Ordinal also barely selects, keeping all four attributes on 76.8\% of pairs. The attribute
decisive learns to discard is \texttt{venue}: the two sources render it incompatibly
(\texttt{"SIGMOD Record"} against \texttt{"The VLDB Journal --- The International Journal
on\ldots"}), so venue disagreement is evidence of formatting, not of non-identity. Scored in
isolation it looks moderately informative, which is why ordinal keeps it.

\subsection{The selection rule matters as much as the elicitation}
\label{sec:rule}

\begin{table}[h]
\centering
\begin{tabular}{lrrrrr}
\toprule
Rule & Attrs/pair & FP & FN & F1 & \texttt{title} only \\
\midrule
cumulative 0.8 (default) & 2.24 & 535 & 10 & 0.8904 &  0.0\% \\
cumulative 0.6           & 1.66 & 434 & 13 & 0.9082 & 34.5\% \\
cumulative 0.5           & 1.27 & 293 & 20 & 0.9337 & 73.2\% \\
cumulative 0.4           & 1.18 & 275 & 26 & 0.9359 & 82.7\% \\
gap (no threshold)       & 1.42 & \textbf{255} & 27 & \textbf{0.9397} & 74.9\% \\
\midrule
Global baseline, fixed \texttt{[title]} & 1.00 & 339 & 26 & 0.9233 & 100.0\% \\
\bottomrule
\end{tabular}
\caption{DBLP--ACM, decisive profiling throughout, varying only the selection rule.}
\end{table}

At the default target of 0.8, pair-adaptive selection \emph{loses} to the global baseline
($p = 1.6\times10^{-13}$); at 0.4 it wins ($66/2$, $p = 1.6\times10^{-17}$). The comparison
between adaptive and global selection turned on an untuned hyperparameter rather than on
the methods.

The best configuration is at neither extreme: showing everything (4.00 attributes) and
showing one fixed attribute (1.00) are both beaten by rules showing roughly 1.2--1.4.
Targets 0.4 and 0.5 are statistically indistinguishable ($55/67$, $p = 0.32$), so this is a
plateau rather than one fortunate value, while the steps above them are individually
significant. The \texttt{gap} rule is the practically important row: it has no threshold, so
nothing about it was chosen by inspecting these results, and it reaches the top of the range
anyway. Its margin over the best hand-tuned target is thin ($50/31$, $p = 0.045$) and should
not be claimed as superiority --- the defensible claim is that it matches a tuned threshold
without requiring one.

\subsection{Why adaptation beats a fixed subset}
\label{sec:mechanism}

The global baseline reduces every record to \texttt{title}, which is also what the adaptive
method chooses for three quarters of pairs. Comparing them isolates adaptation itself,
holding the attribute vocabulary fixed.

\begin{table}[h]
\centering
\begin{tabular}{lrrr}
\toprule
Pair group & Pairs & FP, adaptive & FP, fixed \texttt{[title]} \\
\midrule
Adaptive also chose \texttt{[title]} & 9{,}803 & 118 & 118 \\
Adaptive chose more                  & 3{,}277 & \textbf{137} & 221 \\
\midrule
\quad \texttt{[title, authors, year]} & 2{,}257 & 101 & 146 \\
\quad \texttt{[title, year]}          &    841 &  \textbf{13} &  58 \\
\quad \texttt{[title, authors]}       &    179 &  23 &  17 \\
\bottomrule
\end{tabular}
\caption{DBLP--ACM, all 13{,}080 pairs: decisive under \texttt{gap} against the 200-label
global baseline.}
\end{table}

Where the two select the same attribute they produce identical results --- 118 false
positives each --- which is what the shared cache guarantees and which confirms the
comparison is clean. The entire gain comes from the 3{,}277 pairs where the adaptive method
kept more, cutting false positives from 221 to 137. The clearest single effect is the 841
pairs given \texttt{[title, year]}, where false positives fall from 58 to 13: titles similar
enough to mislead a title-only matcher, with the year settling the question. A global method
cannot express this, because adding \texttt{year} everywhere is close to what the
cumulative-0.8 configuration does, and it costs more elsewhere than it saves here.

\subsection{The matcher is not a fixed quantity}
\label{sec:draws}

Every number reported above is a single call of the matcher over each pair. That is only a
sound way to report a result if repeating the run reproduces it, and it does not.
OpenRouter serves one model slug from several providers that differ in quantization, so two
fresh runs of the same configuration over the same pairs are two different measurements.
Re-running the two headline configurations three times each, with the response cache
disabled so that each run is an independent draw rather than a replay, gives:

\begin{table}[h]
\centering
\begin{tabular}{lrrrrr}
\toprule
Configuration & Draw 1 & Draw 2 & Draw 3 & Mean & Range \\
\midrule
Step 1.5, decisive (\texttt{gap}) & 0.9228 & 0.8886 & 0.9151 & \textbf{0.9088} & 0.0342 \\
No selection                      & 0.7632 & 0.7006 & 0.7047 & 0.7228 & 0.0626 \\
\bottomrule
\end{tabular}
\caption{DBLP--ACM F1 over three independent matcher draws per configuration, all run
26--27 August with \texttt{--disable-matcher-cache}. Same model slug, same prompt hash,
same attribute counts as the runs in Table~1.}
\label{tab:draws}
\end{table}

Three things follow, and they point in different directions.

\textbf{The comparison is robust and larger than first reported.} The two ranges do not
overlap: the worst adaptive draw, 0.8886, exceeds the best no-selection draw, 0.7632. The
separation between the means is 0.186, against the 0.069 that the single draws of Table~1
suggested. The paired tests within each draw agree, giving adaptive selection 1{,}217, 1{,}689
and 1{,}677 uniquely correct pairs against 207, 341 and 224 for no selection
($p < 10^{-170}$ in all three). The finding replicates.

\textbf{The absolute values do not reproduce.} No-selection F1 was 0.8706 in mid-August and
averages 0.7228 a fortnight later, with precision falling from 0.7729 to 0.5674. The cause
is visible in the raw counts: against 2{,}224 true matches the matcher asserted 2{,}867 in
August and 3{,}591--4{,}111 now. It has become markedly more willing to answer Yes, and the
configuration that shows it more fields suffers more for it. Nothing in this pipeline
changed --- the model slug, the prompt hash, the attribute counts and the token counts are
identical across the two periods. Any single-draw F1 in this report should therefore be read
as a measurement of one matcher on one date, not as a property of the method, and a reader
reproducing these runs today should expect the levels to differ while the ordering holds.

\textbf{Condensation halves the variance.} The adaptive configuration varies by 0.034 across
draws and no selection by 0.063. This was not something the method was designed to do, but it
follows from the same mechanism as the accuracy gain: most of what provider variation can
change is a borderline judgement on a field that carries little evidence, and a pair reduced
to 1.42 attributes has fewer such fields to be swayed by. Stability under a drifting
serving stack is a second, independent argument for condensing, and one that matters more in
deployment than in a benchmark.

\subsection{Amazon--Walmart: the wide-schema regime}
\label{sec:amazon}

\begin{table}[h]
\centering
\resizebox{\textwidth}{!}{%
\begin{tabular}{llrrrrrrrr}
\toprule
Method & Rule & Attrs/pair & Tokens & TP & FP & FN & F1 & Precision & Recall \\
\midrule
No selection (\texttt{full}) & --- & 12.00 & 1035 & 1010 & 1018 & 133 & \textbf{0.6370} & 0.4980 & 0.8836 \\
Step 1.5, decisive & \texttt{gap}, $k\!\ge\!1$ & 1.34 & \textbf{248} & 914 & \textbf{919} & 229 & 0.6142 & \textbf{0.4986} & 0.7997 \\
Step 1.5, decisive & \texttt{cumulative} & 2.24 & 286 & 954 & 1273 & 189 & 0.5662 & 0.4284 & 0.8346 \\
Step 1.5, decisive & \texttt{gap}, $k\!\ge\!2$ & 2.39 & 316 & 961 & 1413 & 182 & 0.5465 & 0.4048 & 0.8408 \\
Step 1.5, ordinal & \texttt{gap}, $k\!\ge\!1$ & 2.18 & 382 & 979 & 1494 & 164 & 0.5412 & 0.3956 & 0.8565 \\
Step 1.5, decisive & \texttt{gap}, $k\!=\!5$ & 5.00 & 524 & 934 & 3238 & 209 & 0.3515 & 0.2239 & 0.8171 \\
Step 1, supervised & \multicolumn{9}{l}{did not run: a fixed cut-off selected no attributes} \\
\bottomrule
\end{tabular}}
\caption{Amazon--Walmart, 1{,}143 ground-truth matches, blocking completeness 0.9143.
Metrics are end-to-end, so false negatives include the 98 matches blocking did not retain.
All Step-1.5 rows use the evidence mask of Section~\ref{sec:mask}. The \texttt{cumulative}
row runs with the attribute cap removed, without which the rule cannot reach its target on
this schema at all.}
\end{table}

This is the benchmark with headroom. Baseline precision is 0.4980 here against 0.7729 on
DBLP--ACM and 0.9469 on Fodors--Zagat --- close to half the asserted matches are wrong. It
is also where condensation is worth most: 1{,}035 prompt tokens per pair at the full
representation, against 306 on DBLP--ACM.

\textbf{Adaptive selection matches the baseline's per-pair accuracy at a quarter of the
prompt.} Over the 88{,}296 shared candidate pairs the baseline is uniquely correct on 804
and decisive selection on 807, a paired difference of three pairs in nearly ninety thousand
($p = 0.96$). The two are, on this evidence, equally accurate matchers, and one of them
reads 248 tokens per pair rather than 1{,}035.

F1 nonetheless favours the baseline, 0.6370 to 0.6142, and the two statements are
consistent rather than contradictory: equal error \emph{counts} need not mean equal F1 when
the errors fall differently. Precision is identical to four decimal places (0.4986 against
0.4980); the entire F1 gap is recall, 0.7997 against 0.8836. With only 1{,}143 matches among
88{,}296 pairs, a false negative costs recall far more than a false positive costs precision,
so an even trade of one for the other still moves F1. Which of the two numbers matters is an
application question, not a methodological one: a deployment paying per token and rescoring
survivors should read the paired test, while one that must not miss matches should read F1.

\textbf{No fixed rule recovers the difference; every one tried makes it worse.} The recall
gap sits entirely in pairs condensed to a single attribute, which suggested raising the
floor. Doing so does raise recall --- false negatives fall from 229 to 182 --- but precision
collapses from 0.4986 to 0.4048 and F1 with it. Splitting by what the unconstrained selector
had chosen isolates the cause exactly: on the 27{,}402 pairs it had already given two or more
attributes nothing changes, false positives 150 against 150, while on the 60{,}894 pairs it
had condensed to one, forcing a second attribute buys 47 true positives at the price of 494
false ones. The attribute forced in is \texttt{brand} for 33{,}823 of them, and two different
products from one manufacturer agree on it; the floor injects apparent corroboration where
the selector had correctly found none. A second rule reaching a similar width by an unrelated
mechanism --- \texttt{cumulative} with its cap removed, 2.24 attributes --- reproduces the
same collapse (precision 0.4284), so this is a property of the two-attribute region on this
schema rather than an artefact of the floor. Both are significant against the unconstrained
selector: $p = 3.5\times10^{-27}$ and $p = 1.3\times10^{-19}$.

\textbf{The interval between those two widths is worse than either end of it.} Forcing five
attributes per pair --- requested as six, capped at five --- collapses precision to 0.2239
and F1 to 0.3515, far below both the 1.34-attribute configuration and no selection. It loses
to no selection on 2{,}882 pairs against 586, and to the unconstrained selector on 3{,}028
against 729 ($p$ below double precision in both). Precision as a function of width is thus
not merely non-monotone but sharply valleyed: 0.4986 at 1.34 attributes, 0.4048 at 2.39,
0.2239 at 5, and 0.4980 again at 12.

Counting what the five-attribute prompts actually contain explains the shape. The subset is
led by \texttt{title} (100\% of pairs), \texttt{brand} (90.6\%), \texttt{modelno} (71.1\%)
and \texttt{longdescr} (61.5\%) --- descriptive text, on which distinct products routinely
agree. The fields that let a matcher refuse appear rarely: \texttt{price} in 26.1\% of pairs
and \texttt{shipweight} in 0.7\%, because importance ranks them low and the cut removes them
first. Five attributes therefore supply abundant corroboration and almost no grounds for
contradiction, which is precisely the composition that manufactures false positives. At
twelve attributes the numeric fields return and precision recovers.

The practical consequence is that this schema offers no cheap middle. Either condense hard,
to 1.34 attributes for a quarter of the tokens and no loss of per-pair accuracy, or show
everything. Every intermediate width tested is worse than both.

\textbf{Elicitation format replicates across datasets; the schema effect does not.} Decisive
beats ordinal here as it did on DBLP--ACM, 1{,}230 pairs to 806 ($p = 5.0\times10^{-21}$).
That ordering is the one finding in this section that carries over unchanged.

The supervised baseline could not be run here. Its importance scores are normalised to sum
to one and it retained attributes above a fixed 0.3 cut-off, which on a 12-attribute schema
no attribute reaches --- the mean share is $1/12$. It selected the empty subset and aborted.
This is the same scale-free defect as the cumulative rule of Section~\ref{sec:rule}, in a
different component.

One structural finding came out of preparing this dataset. The \texttt{original\_id} field
is a per-table row counter running $1\ldots22{,}074$ and $1\ldots2{,}554$, sharing no key
across the tables. Excluding it raised pair completeness from 0.9099 to 0.9143, recovering
five ground-truth pairs. Because each column is encoded independently and the column vectors
are concatenated, that field occupied a full block of every record embedding and displaced
real matches out of the candidate set. No attribute selection could have recovered them,
since Step 1.5 does not re-block --- the repair had to happen before embedding.

\section{Discussion}

\textbf{Pair-adaptive selection outperforms both no selection and a labelled global
baseline --- on DBLP--ACM.} There it improves F1 from 0.8706 to 0.9397 while showing the
matcher 65\% fewer attributes ($p = 5.2\times10^{-50}$), and beats a baseline trained on 400
labelled pairs over the identical candidate set ($p = 1.0\times10^{-12}$) while using none.
The qualifier is not decorative: on Amazon--Walmart the same configuration only draws with
no selection, and draws on per-pair accuracy while still trailing on F1
(Section~\ref{sec:amazon}). The improvement is a result about one dataset; what generalises
is weaker and is stated separately below.

\textbf{How importance is elicited decides whether adaptation helps at all.} The ordinal
formulation, on the same pipeline and rule, is significantly worse than no selection.
Scoring each attribute in isolation produces a nearly flat importance vector that no rule
can act on; asking for a sufficient subset forces the comparison the task requires. This is
the most transferable finding here, and it is a property of the elicitation rather than of
the dataset.

\textbf{The rule consuming the importance vector is as consequential as the vector.} The
same profiling output, read by the default rule, loses to the global baseline; read by
\texttt{gap}, it beats it. Any claim about pair-adaptive selection that does not state its
selection rule is under-specified. The default was not merely suboptimal but, on wide
schemas, structurally incapable of condensing at all.

\textbf{The benefit is conditional on the schema, and the condition is now measured.} On
Fodors--Zagat, saturated, no method differs from any other. On DBLP--ACM, four attributes
with \texttt{title} close to a unique key, adaptation wins. On Amazon--Walmart, twelve
attributes with no such key, it draws: condensing to 1.34 attributes holds precision exactly
and costs recall. The pattern across the three is that aggressive condensation wins outright
where one or two attributes can carry the decision, and degrades to a cost-saving with no
accuracy gain where the evidence is spread across many. Nothing in the method currently
detects which regime it is in --- though the importance vector it already computes contains
the information, since the regimes differ precisely in how concentrated that vector is.

\textbf{The matcher is a larger source of variation than most of the effects studied here.}
Repeating one configuration three times moves DBLP--ACM F1 by 0.03 for adaptive selection and
0.06 for no selection, and two weeks of provider drift moved the no-selection level by 0.15
with nothing in the pipeline changed. Several differences reported on Amazon--Walmart are
smaller than that. The comparisons remain sound, because each is paired within a single draw,
but the levels are properties of a matcher on a date. A useful side effect is that
condensation roughly halves this variance (Section~\ref{sec:draws}), which is an argument for
it that has nothing to do with accuracy.

\textbf{Adaptivity earns its keep most clearly where it wins least.} On the schema where
condensation does not improve F1, every fixed rule tried in its place is significantly
\emph{worse} than letting the model choose per pair: a floor of two attributes
($p = 3.5\times10^{-27}$), and a cumulative target reaching a comparable width by an
unrelated mechanism ($p = 1.3\times10^{-19}$). The selector's most aggressive decisions ---
the 60{,}894 pairs it cut to a single attribute --- are the ones a fixed rule most wants to
overrule, and overruling them costs 494 false positives to buy 47 true ones. This is the
sharpest evidence in the report that per-pair selection is doing something a global subset
cannot, and it comes from the dataset where the method does not win.

\section{Limitations}

\textbf{The supervised global baseline is missing on Amazon--Walmart.} Its fixed cut-off
selected nothing on a 12-attribute schema and the run aborted, so no \emph{labelled} global
method is compared there. The unlabelled global comparison it was meant to supply has since
been made another way: the fixed-floor and cumulative configurations of
Section~\ref{sec:amazon} apply a near-constant subset width to every pair, and both are
significantly worse than per-pair selection. What remains untested on this schema is whether
a global subset chosen \emph{with labels} would do better than either.

\textbf{The selection rule was chosen on DBLP--ACM; its confirmation is outstanding.} The
sweep in Section~\ref{sec:rule} was read on the same pairs it reports over, so the
cumulative targets in it are tuned values and none should be quoted as the method's
performance. The rule adopted, \texttt{gap}, has no threshold and so contributes no tuned
quantity, but preferring it over \texttt{cumulative} was itself a choice made after seeing
those results.

The design handles this by role separation: DBLP--ACM is the \emph{selection} dataset and
Amazon--Walmart the \emph{held-out} one. \texttt{gap} was fixed before any Amazon--Walmart
selection run started and has not been revised since, so the rule itself contributes no
quantity tuned on held-out data.

\textbf{That separation has since been broken once, and the Amazon--Walmart figures should
be read accordingly.} The evidence mask of Section~\ref{sec:mask} was found by inspecting
Amazon--Walmart failures, and it raises decisive selection there from 0.5959 to 0.6142. The
current Amazon--Walmart numbers are therefore no longer a clean out-of-sample estimate:
one component of the method was derived from the test set. Two considerations limit the
damage without removing it. The mask is a correctness fix with an argument independent of
any measurement --- an attribute blank on both sides cannot inform a decision, whatever the
data says --- rather than a parameter fitted to improve a score. And it is inert on the
selection dataset, changing 5 of 13{,}080 pairs and leaving DBLP--ACM F1 unchanged, so it
cannot have been tuned to a result there either. Restoring the guarantee needs a third
dataset that has been used for nothing else; until then the honest reading of the
Amazon--Walmart column is a strong estimate with a known dependency on the data it is
computed over.

\textbf{Most configurations are still a single matcher draw.} Section~\ref{sec:draws}
repeats the two DBLP--ACM headline rows three times each and finds the ordering stable and
the levels not. Every other row in this report --- the whole of Fodors--Zagat, every
Amazon--Walmart configuration, the rule sweep, and the four DBLP--ACM rows not repeated ---
remains one draw, and the drift measured on DBLP--ACM gives no reason to assume the others
are steadier. The Amazon--Walmart comparisons are the ones to treat most carefully, since
several of the differences there are smaller than the 0.06 spread the repeated DBLP--ACM
runs exhibited. The paired design makes each comparison valid within its own draw; it does
not make the levels reproducible. The large DBLP--ACM effects are unlikely to reverse; the
\texttt{gap}-versus-tuned-threshold comparison, with 81 discordant pairs and $p = 0.045$,
could.

\textbf{The supervised baseline is not an equal competitor.} It consumes 400 labelled pairs
while the adaptive method consumes none, so the comparison measures method and label budget
together. Its selected subset also moves with the budget --- \texttt{[title, authors]} at 10
and 50 labels, \texttt{[title]} at 200 --- so no single sample size characterises it.

\textbf{The feature vector is weaker than it appears.} Of the six per-attribute features,
\texttt{semantic\_similarity} is disabled and constant at zero and \texttt{is\_numeric} is
constant per column on both datasets reported here; on the widest schema
\texttt{exact\_match} fires on 0.06\% of cells. Effectively three features carry signal. The
vector also has no notion of value rarity, the standard discriminative signal in entity
resolution: matching on a rare title token and on a common one are represented identically.

\textbf{One sampler.} All Step 1.5 runs use the \texttt{diversity} profile sampler, held
fixed so the other variables could be read. No claim about sampler choice is made.

\end{document}
